% GENERATED FILE. Edit canonical lessons, then run npm run generate:books.
% canonical-source-hash: fnv1a64:74c5914d64e810af
% canonical-lessons: GU-C313-aphva, GU-C313-phelavvu, GU-C313-satya, GU-C313-mahiti, GU-C313-khabar

\chapter{A Rumour Spreads}
\label{ch:gu-a-rumour-spreads}
\hlchaptermodality{313}

\begin{chapteropening}
\textbf{By the end of this chapter:} \emph{I can tell a rumour from the truth.}

Complete the last lesson of chapter 313: I can tell a rumour from the truth.
\end{chapteropening}

\section[\texorpdfstring{\gu{અફવા} (aphvā)}{અફવા (aphvā)}]{\gu{અફવા} (aphvā) — a rumour}

\label{lesson:GU-C313-aphva}



\begin{quote}
Before the new one: say the Gujarati for a report, then the Gujarati for to describe.
\end{quote}

\subsection*{You'll want to know: \gu{અફવા}}
\textbf{\gu{અફવા}} — \emph{aphvā} — \textquotedblleft{}a rumour\textquotedblright{}.

\begin{quote}
\gu{ગામમાં} \gu{અફવા} \gu{ફેલાઈ} \gu{કે} \gu{ડાકુ} \gu{આવ્યા} \gu{છે}.
\end{quote}

\begin{quote}
\emph{gāmmā̃ aphvā phelāī ke ḍāku āvyā che.}
\end{quote}

\begin{quote}
A rumour went round the village that bandits had come.
\end{quote}

A feminine noun.

\begin{cousinweb}
\textbf{\gu{અફવા}} (\emph{aphvā}) is usually traced, through Persian, to Arabic \emph{afwāh}, \textquotedblleft{}mouths\textquotedblright{}: what passes from mouth to mouth. It \textbf{\gu{ફેલાય} \gu{છે}}, spreads.
\end{cousinweb}

\begin{grammarlens}[title={What you've built}]
One more word for telling what happened.
\end{grammarlens}

\subsection*{Your turn}
\begin{itemize}
\raggedright
  \item \emph{Say it:} \emph{aphvā}
  \item \emph{Say it:} \emph{aphvā}, once more
  \item \emph{Say it:} \emph{gāmmā̃ aphvā phelāī ke ḍāku āvyā che.} — A rumour went round the village that bandits had come.
  \item \emph{Recall:} say the Gujarati for a report, then the Gujarati for to describe, then say \emph{aphvā} again
\end{itemize}

\begin{tcolorbox}[breakable,colback=teal!4,colframe=teal!35!black,title={Before you move on}]
What does \gu{અફવા} mean? (\textquotedblleft{}A rumour\textquotedblright{}.) Say it once more.
\end{tcolorbox}
\section[\texorpdfstring{\gu{ફેલાવવું} (phelāvvũ)}{ફેલાવવું (phelāvvũ)}]{\gu{ફેલાવવું} (phelāvvũ) — to spread (something)}

\label{lesson:GU-C313-phelavvu}



\begin{quote}
Before the new one: say the Gujarati for to describe, then the Gujarati for a rumour.
\end{quote}

\subsection*{You'll want to know: \gu{ફેલાવવું}}
\textbf{\gu{ફેલાવવું}} — \emph{phelāvvũ} — \textquotedblleft{}to spread (something)\textquotedblright{}.

\begin{quote}
\gu{કોઈએ} \gu{જાણી} \gu{જોઈને} \gu{આ} \gu{અફવા} \gu{ફેલાવી}.
\end{quote}

\begin{quote}
\emph{koīe jāṇī joīne ā aphvā phelāvī.}
\end{quote}

\begin{quote}
Someone spread this rumour on purpose.
\end{quote}

Past, with \gu{તેણે} (he or she …): \gu{તેણે} \gu{અફવા} \gu{ફેલાવી}, \emph{teṇe aphvā phelāvī}.

\begin{cousinweb}
\textbf{\gu{ફેલાવવું}} (\emph{phelāvvũ}) is the causative of \textbf{\gu{ફેલાવું}}, to spread, which goes back to Sanskrit \emph{prathita}, \textquotedblleft{}spread out\textquotedblright{} (CDIAL 8651): to make something spread. \textbf{\gu{જાણી} \gu{જોઈને}}, knowing and seeing, is on purpose.
\end{cousinweb}

\begin{grammarlens}[title={What you've built}]
One more word for telling what happened.
\end{grammarlens}

\subsection*{Your turn}
\begin{itemize}
\raggedright
  \item \emph{Say it:} \emph{phelāvvũ}
  \item \emph{Say it:} \emph{phelāvvũ}, once more
  \item \emph{Say it:} \emph{koīe jāṇī joīne ā aphvā phelāvī.} — Someone spread this rumour on purpose.
  \item \emph{Recall:} say the Gujarati for to describe, then the Gujarati for a rumour, then say \emph{phelāvvũ} again
\end{itemize}

\begin{tcolorbox}[breakable,colback=teal!4,colframe=teal!35!black,title={Before you move on}]
What does \gu{ફેલાવવું} mean? (\textquotedblleft{}To spread (something)\textquotedblright{}.) Say it once more.
\end{tcolorbox}
\section[\texorpdfstring{\gu{સત્ય} (satya)}{સત્ય (satya)}]{\gu{સત્ય} (satya) — the truth}

\label{lesson:GU-C313-satya}



\begin{quote}
Before the new one: say the Gujarati for a rumour, then the Gujarati for to spread (something).
\end{quote}

\subsection*{You'll want to know: \gu{સત્ય}}
\textbf{\gu{સત્ય}} — \emph{satya} — \textquotedblleft{}the truth\textquotedblright{}.

\begin{quote}
\gu{પોલીસે} \gu{લોકોને} \gu{સત્ય} \gu{જણાવ્યું}.
\end{quote}

\begin{quote}
\emph{polīse lokone satya jaṇāvyũ.}
\end{quote}

\begin{quote}
The police told people the truth.
\end{quote}

A neuter noun.

\begin{cousinweb}
\textbf{\gu{સત્ય}} (\emph{satya}) is a tatsama, Sanskrit \emph{satya}, \textquotedblleft{}true, real\textquotedblright{} (Monier-Williams). The noun is neuter; \textbf{\gu{સાચું}} is the everyday adjective, true.
\end{cousinweb}

\begin{grammarlens}[title={What you've built}]
One more word for telling what happened.
\end{grammarlens}

\subsection*{Your turn}
\begin{itemize}
\raggedright
  \item \emph{Say it:} \emph{satya}
  \item \emph{Say it:} \emph{satya}, once more
  \item \emph{Say it:} \emph{polīse lokone satya jaṇāvyũ.} — The police told people the truth.
  \item \emph{Recall:} say the Gujarati for a rumour, then the Gujarati for to spread (something), then say \emph{satya} again
\end{itemize}

\begin{tcolorbox}[breakable,colback=teal!4,colframe=teal!35!black,title={Before you move on}]
What does \gu{સત્ય} mean? (\textquotedblleft{}The truth\textquotedblright{}.) Say it once more.
\end{tcolorbox}
\section[\texorpdfstring{\gu{માહિતી} (māhitī)}{માહિતી (māhitī)}]{\gu{માહિતી} (māhitī) — information}

\label{lesson:GU-C313-mahiti}



\begin{quote}
Before the new one: say the Gujarati for to spread (something), then the Gujarati for the truth.
\end{quote}

\subsection*{You'll want to know: \gu{માહિતી}}
\textbf{\gu{માહિતી}} — \emph{māhitī} — \textquotedblleft{}information\textquotedblright{}.

\begin{quote}
\gu{આ} \gu{માહિતી} \gu{ખોટી} \gu{હતી}.
\end{quote}

\begin{quote}
\emph{ā māhitī khoṭī hatī.}
\end{quote}

\begin{quote}
This information was wrong.
\end{quote}

A feminine noun.

\begin{cousinweb}
\textbf{\gu{માહિતી}} (\emph{māhitī}) is usually traced, through Persian, to Arabic \emph{māhiyya}, \textquotedblleft{}what a thing is, its nature\textquotedblright{}: the facts of a matter. It is feminine and does not take a plural: \textbf{\gu{ખોટી} \gu{માહિતી}}, wrong information.
\end{cousinweb}

\begin{grammarlens}[title={What you've built}]
One more word for telling what happened.
\end{grammarlens}

\subsection*{Your turn}
\begin{itemize}
\raggedright
  \item \emph{Say it:} \emph{māhitī}
  \item \emph{Say it:} \emph{māhitī}, once more
  \item \emph{Say it:} \emph{ā māhitī khoṭī hatī.} — This information was wrong.
  \item \emph{Recall:} say the Gujarati for to spread (something), then the Gujarati for the truth, then say \emph{māhitī} again
\end{itemize}

\begin{tcolorbox}[breakable,colback=teal!4,colframe=teal!35!black,title={Before you move on}]
What does \gu{માહિતી} mean? (\textquotedblleft{}Information\textquotedblright{}.) Say it once more.
\end{tcolorbox}
\section[\texorpdfstring{\gu{ખબર} (khabar)}{ખબર (khabar)}]{\gu{ખબર} (khabar) — knowledge (of something); word}

\label{lesson:GU-C313-khabar}



\begin{quote}
Before the new one: say the Gujarati for the truth, then the Gujarati for information.
\end{quote}

\subsection*{You'll want to know: \gu{ખબર}}
\textbf{\gu{ખબર}} — \emph{khabar} — \textquotedblleft{}knowledge (of something); word\textquotedblright{}.

\begin{quote}
\gu{મીરાને} \gu{ચોરીની} \gu{ખબર} \gu{સવારે} \gu{પડી}.
\end{quote}

\begin{quote}
\emph{Mīrāne corīnī khabar savāre paḍī.}
\end{quote}

\begin{quote}
Mira found out about the theft in the morning.
\end{quote}

A feminine noun.

\begin{cousinweb}
\textbf{\gu{ખબર}} (\emph{khabar}) is usually traced, through Persian, to Arabic \emph{khabar}, \textquotedblleft{}news, a report\textquotedblright{}. In Gujarati it is knowing about something: \textbf{\gu{ખબર} \gu{પડી}}, came to know; \textbf{\gu{મને} \gu{ખબર} \gu{છે}}, I know.
\end{cousinweb}

\begin{grammarlens}[title={What you've built}]
One more word for telling what happened.
\end{grammarlens}

\subsection*{Your turn}
\begin{itemize}
\raggedright
  \item \emph{Say it:} \emph{khabar}
  \item \emph{Say it:} \emph{khabar}, once more
  \item \emph{Say it:} \emph{Mīrāne corīnī khabar savāre paḍī.} — Mira found out about the theft in the morning.
  \item \emph{Recall:} say the Gujarati for the truth, then the Gujarati for information, then say \emph{khabar} again
\end{itemize}

\begin{tcolorbox}[breakable,colback=teal!4,colframe=teal!35!black,title={Before you move on}]
What does \gu{ખબર} mean? (\textquotedblleft{}Knowledge (of something); word\textquotedblright{}.) Say it once more.
\end{tcolorbox}
